\documentclass[11pt]{article}
% common.tex — minimal noise setup
% ----------------------------

% Silence package messages before anything else
\RequirePackage{silence}

% Completely hide "info" messages from LaTeX packages
\WarningFilter*{latex}{}
\WarningFilter*{hyperref}{}
\WarningFilter*{pgf}{}
\WarningFilter*{tikz}{}

% Optionally: hide some hyperref PDF string warnings that are common but harmless
\WarningFilter{hyperref}{Token not allowed in a PDF string}

% ----------------------------
% Now load the usual packages
\usepackage[utf8]{inputenc}
\usepackage[T1]{fontenc}
\usepackage{amsmath, amssymb, amsthm, amsfonts}
\usepackage{geometry}
\usepackage{hyperref}
\usepackage{graphicx}
\usepackage{physics}
\usepackage{booktabs}
\usepackage{listings}
\usepackage{bookmark}
\usepackage{amscd}
\usepackage{tikz-cd}
\usepackage{mathtools}
\usepackage{xspace}
\usepackage[backend=biber, style=numeric]{biblatex}
\usepackage{glossaries}
\usepackage{tcolorbox}

% ----------------------------
% Theorems, macros
\newtheorem{lemma}{Lemma}
\newcommand{\IaMe}{\(IaM^e\)\xspace}
\newcommand{\pdflink}[1]{}

% ----------------------------
% Page layout
\geometry{letterpaper, margin=1in}
\setlength{\parindent}{0pt}
\setlength{\parskip}{1em}

% Hyperref settings
\hypersetup{
    colorlinks=true,
    linkcolor=black,
    citecolor=black,
    filecolor=magenta,
    urlcolor=blue
}

% TOC formatting
\usepackage{tocloft}
\setlength{\cftbeforechapskip}{1.0em}
\setlength{\cftbeforesecskip}{0pt}
\setlength{\cftbeforesubsecskip}{0pt}
\renewcommand{\cftchapleader}{\cftdotfill{\cftdotsep}}

% Glossary
\defglsentryfmt{\ifglsused{\glslabel}{\glsgenentryfmt}{\textit{\glsgenentryfmt}}}

% Author
\author{$IaM^e$}


\input{../glossary.tex}
\addbibresource{../references.bib}

\title{Quantum Cosmology: \\ Spectral Complexity and the Measure of Observers in a Static Universe}

\date{2025}

\begin{document}

\maketitle

\begin{abstract}
  We propose a purely information-theoretic formulation of quantum cosmology that operates entirely within the timeless Wheeler-DeWitt framework.
  Instead of performing a Wick rotation to Euclidean time, we introduce \textit{Spectral Complexity} \(C_{\text{spec}}\) as the fundamental cost functional that determines the measure over static observer wavefunctions. 

  Observers are modeled as relational wavefunctions \(\psi_o(\lambda)\) entangled with a clock variable \(\lambda\) (following the Page-Wootters construction).
  The probability of an observer is weighted by \(P(\psi_o) \propto 2^{-C_{\text{spec}}(\psi_o)}\), where \(C_{\text{spec}}\) counts the effective number of Fourier modes required to describe \(\psi_o(\lambda)\).
  Low-spectral-complexity observers, which perceive smooth and law-like universes, dominate the measure exponentially.

  In a minisuperspace implementation with a scale factor and relational clock, we demonstrate numerically that minimal-\(C_{\text{spec}}\) observer states naturally select classical, expanding FLRW-like trajectories without invoking Euclidean path integrals.
  This framework provides a justification for the emergence of classical spacetime and the principle of least action as a consequence of minimal description length.
  Chaotic histories are suppressed by high spectral complexity, offering a novel resolution to longstanding issues in quantum cosmology.
\end{abstract}

\maketitle

\section{Introduction}

The canonical quantization of general relativity leads to the Wheeler-DeWitt equation \cite{DeWitt1967}:
\begin{equation}
\hat{\mathcal{H}} \Psi[g_{ij}, \phi] = 0,
\end{equation}
where \(\Psi\) is the wavefunction of the universe on superspace.
This equation is timeless: there is no external time parameter.
The ``Problem of Time'' \cite{isham1993canonical} arises because all physical predictions must be extracted from a single, static superposition.

The standard approach to recovering dynamics and probabilities relies on the Hartle-Hawking no-boundary proposal \cite{hartle1983wave}, which employs a Wick rotation to Euclidean time and weights histories by \(e^{-I_{\text{Eucl}}}\).
While successful in many minisuperspace models, this method faces conceptual difficulties: the physical meaning of imaginary time, the conformal factor problem, and the justification of the Euclidean measure itself remain debated.

In this work we propose \textit{Spectral Quantum Cosmology} (SQC).
We remain strictly within the static Wheeler-DeWitt framework and introduce an information-theoretic selection principle based on \textit{Spectral Complexity}.
Rather than weighting entire four-geometries via Euclidean action, we weight possible \textit{observer wavefunctions} directly.
Observers with simpler (more compressible) internal descriptions dominate the measure.
This naturally selects observers who experience classical, smooth, law-like spacetimes.

\section{Relational Observers in a Timeless Universe}

Following Page and Wootters \cite{page1983clock,wootters1984time}, time emerges relationally.
The universal wavefunction \(\Psi\) is static and contains entangled subsystems: a clock degree of freedom and the rest of the universe.
An observer corresponds to a wavefunction \(\psi_o\) that is entangled with a suitable clock variable \(\lambda\) (which may be the scale factor, a scalar field value, or an internal memory register).

The conditional probability that the rest of the universe is in state \(|\chi\rangle\) when the observer's clock reads \(\lambda\) is
\begin{equation}
P(\chi | \lambda) = \frac{|\langle \chi | \psi_o(\lambda) \rangle|^2}{\langle \psi_o(\lambda) | \psi_o(\lambda) \rangle}.
\end{equation}
Dynamics and history are recovered from correlations within this timeless structure.
The central question becomes: \textit{which observer wavefunctions \(\psi_o(\lambda)\) have the highest measure in the universal superposition?}

\section{Spectral Complexity as the Measure}

We define the \textit{Spectral Complexity} of an observer wavefunction \(\psi_o(\lambda)\) as the minimal information cost of representing it in the frequency domain.
Let
\begin{equation}
\psi_o(\lambda) = \sum_{k=1}^{M} c_k \, e^{i \omega_k \lambda}
\end{equation}
be the Fourier decomposition with respect to the relational clock \(\lambda\), where only modes with significant amplitude \(|c_k|\) are retained.

The spectral complexity is then
\begin{equation}
C_{\text{spec}}(\psi_o) \approx M \cdot \bigl( b + \log_2(\omega_{\max}/\omega_{\min}) \bigr) + \alpha \, \mathcal{D},
\end{equation}
where \(b\) is the bit precision of the coefficients, \(\omega_{\max}\) is the highest relevant frequency, and \(\mathcal{D}\) is a phase-incoherence penalty (e.g., integrated Fubini-Study distance between neighboring \(\lambda\) slices) that favors continuity.

According to Solomonoff induction and algorithmic probability \cite{Solomonoff1964}, the natural prior weight for any observer is
\begin{equation}
P(\psi_o) \propto 2^{-C_{\text{spec}}(\psi_o)} = e^{-(\ln 2) \, C_{\text{spec}}(\psi_o)}.
\end{equation}
Observers whose internal wavefunctions require many Fourier modes (chaotic or rapidly oscillating) are exponentially suppressed.
Conversely, observers who experience smooth, slowly varying, law-like evolution dominate.

This replaces the Euclidean suppression \(e^{-I_{\text{Eucl}}}\) with a native, real-valued algorithmic measure.

\section{Minisuperspace Implementation}

We implement this idea in a closed FLRW minisuperspace with scale factor \(a\) (or \(\alpha = \ln a\)) and a relational clock \(\lambda\) (taken here as \(\alpha\) itself for simplicity, or a homogeneous scalar field).

The discrete configuration space is a one-dimensional lattice over \(\alpha_i\), \(i = 1 \dots N\).
The observer wavefunction \(\psi_o(\alpha_i)\) is represented as a complex vector on this lattice.
We variationally minimize \(C_{\text{spec}}\) subject to normalization, effectively solving the algorithmic constraint
\begin{equation}
\bigl( C_{\text{spec}} - R_{\text{limit}} \bigr) \psi_o \approx 0,
\end{equation}
where \(R_{\text{limit}}\) reflects the finite information capacity of the universe.

Numerical optimization (via gradient descent on Fourier coefficients or direct search) yields dominant low-\(C_{\text{spec}}\) states.
These states concentrate along smooth, classically expanding trajectories.
Chaotic or bouncing solutions near singularities require high-frequency content and are heavily suppressed.

\section{Numerical Results (Proof of Concept)}

Preliminary simulations show strong correlation between low spectral complexity and classical least-action paths.
Dominant observer wavefunctions \(\psi_o(\alpha)\) exhibit few Fourier modes and peak near the classical Friedmann solutions for radiation- and matter-dominated epochs.
As lattice resolution increases, the extracted conditional trajectories converge toward the expected classical expansion history.

The effective ``action'' derived from \(-\ln P(\psi_o)\) reproduces the qualitative features of the Euclidean action in standard quantum cosmology, supporting the identification \(I_{\text{Eucl}} \approx \gamma \, C_{\text{spec}}\) in the semiclassical limit.

(Figures showing \(|\psi_o(\alpha)|^2\) for minimal-\(C_{\text{spec}}\) states, comparison with classical solutions, and scaling with resolution will appear in the full version.)

\section{Discussion}

This framework offers several conceptual advantages:

1. \textbf{Timelessness is fundamental.} No Wick rotation or external time is required.
2. \textbf{Native real probabilities.} The measure \(2^{-C_{\text{spec}}}\) is positive-definite and algorithmically motivated.
3. \textbf{Observer selection.} The dominance of low-spectral-complexity observers explains why we experience a smooth, classical universe governed by simple laws.
4. \textbf{Emergence of least action.} The principle of least action arises as the macroscopic manifestation of minimal description length.

Future work will extend the model to inhomogeneous perturbations, investigate the emergence of locality and 3D geometry from spectral compressibility, and study the relational dynamics of entangled observers.



\section{Conclusion}

By weighting observer wavefunctions according to spectral complexity within the static Wheeler-DeWitt universe, we obtain a fully information-theoretic quantum cosmology.
Typical observers, selected by minimal spectral complexity, inhabit smooth classical spacetimes.
This approach bridges quantum gravity, algorithmic information theory, and relational quantum mechanics, suggesting that the universe we observe is the one that is simplest to describe from the inside.


\end{document}
